\chapter{Complete search}\label{ch:search}
\index{complete search}

The cheapest algorithm in the book: walk through every candidate answer and
keep the ones that work. It sounds naive and it is --- and for small enough
search spaces it is also optimal, since the fastest correct program is the one
you finish writing. This chapter teaches you the three big enumeration
idioms --- subsets, permutations, backtracking --- and, more important, the
habit of measuring the search space before committing to it. After this
chapter you will be able to look at a constraint like $n \le 20$, recognise it
as an invitation rather than a warning, and enumerate correctly the first
time.

\section{The exhaustive idea}\label{sec:exhaustive}
\index{brute force|see{complete search}}

A \emph{complete search}\index{complete search!definition} enumerates every
element of the \emph{search space}\index{search space} --- the set of
candidates the answer could be --- and tests each one. Correctness is free:
if the space really contains every candidate, the best survivor is the answer.
The only question that matters is cost, and the cost has one formula:

\[
\text{time} \;=\; (\text{size of search space}) \times (\text{work per candidate}).
\]

Both factors are yours to shrink. The first is fixed by the problem, the
second by your code; keep the per-candidate work small and the formula stays
friendly surprisingly far up. How far? Figure~\ref{fig:ladder} puts the three
sizes you will meet most often next to the Chapter~1 budget.

\begin{figure}[ht]\centering
% fig09-ladder.tex -- Figure 2.4: the size ladder n^2, 2^n, n! vs 1-second
% budget, with the measured enumeration timings of code/size_ladder.cpp.
% Curves given as coordinate lists: factorial() is not a pgfmath function.
\begin{tikzpicture}
\begin{loglogaxis}[
  width=0.86\linewidth, height=6.6cm,
  xlabel={input size $n$}, ylabel={search-space size},
  xmin=5, xmax=40, ymin=1e4, ymax=1e26,
  legend style={at={(0.03,0.97)}, anchor=north west,
                font=\footnotesize, fill=white},
  grid=major, grid style={gray!18},
  major tick style={draw=none},
  axis line style={slate},
  label style={font=\small},
]
\addplot[thick, ink]    coordinates {(5,25) (10,100) (15,225) (20,400)
  (25,625) (30,900) (35,1225) (40,1600)};
\addlegendentry{$n^2$ (for contrast)}
\addplot[thick, accent] coordinates {(5,32) (10,1024) (15,32768) (20,1048576)
  (25,33554432) (27,134217728) (30,1073741824) (33,8589934592)
  (36,6.87e10) (40,1.1e12)};
\addlegendentry{$2^n$}
\addplot[thick, amber]   coordinates {(5,120) (6,720) (7,5040) (8,40320)
  (9,362880) (10,3628800) (11,3.99e7) (12,4.79e8) (13,6.23e9)
  (14,8.72e10) (15,1.31e12) (16,2.09e13) (17,3.56e14) (18,6.40e15)
  (19,1.22e17) (20,2.43e18)};
\addlegendentry{$n!$}
\addplot[thick, amber, dashed] coordinates {(5,1e8) (40,1e8)};
\addlegendentry{1-second budget ($10^8$)}
% measured times, annotated at their n (subsets/sec ~ constant)
\node[font=\scriptsize, text=slate, anchor=west, align=flush left]
  at (axis cs: 6.2, 3e13) {measured, one op per subset (this machine):\\
   $n=20$: 0.24\,ms \quad $n=27$: 35.4\,ms \quad $n=30$: 286\,ms};
% feasibility bands
\draw[slate, dotted, thick] (axis cs:20,1e4) -- (axis cs:20,1e26)
  node[pos=0.93, anchor=south, rotate=90, font=\scriptsize, text=slate]
  {$2^n$ fits to $n \approx 27$};
\draw[slate, dotted, thick] (axis cs:11,1e4) -- (axis cs:11,1e26)
  node[pos=0.88, anchor=south, rotate=90, font=\scriptsize, text=slate]
  {$n!$ fits to $n \approx 11$};
\end{loglogaxis}
\end{tikzpicture}

\caption{The size ladder. Notice where each curve crosses the amber budget
line: $n^2$ is comfortable to $n \approx 10^4$, $2^n$ to $n \approx 27$, and
$n!$ only to $n \approx 11$. The measured points are real: enumerating $2^{30}$
subsets with one cheap operation each took $286$\,ms on this machine --- the
shape of the curve, not your machine's speed, is the lesson.}
\label{fig:ladder}
\end{figure}

The listing behind those measured points is \texttt{code/size\_ladder.cpp};
its timing loop touches each subset once. Real output:

\begin{lstlisting}[numbers=none]
n=15  subsets=32768      time=0.007417 ms
n=20  subsets=1048576    time=0.235667 ms
n=24  subsets=16777216   time=3.76142 ms
n=27  subsets=134217728  time=35.3861 ms
n=29  subsets=536870912  time=143.94 ms
n=30  subsets=1073741824 time=285.851 ms
\end{lstlisting}

Read the doubling, not the absolute numbers: each extra element multiplies
the work by exactly two, because the space itself doubles. That is the whole
shape of exponential growth --- Section~\ref{sec:subsets} is about making
that shape small enough to afford.

\begin{remark}[When $n \le 20$ is a feature]\label{rem:smalln}
\index{constraint!as invitation}
Contest setters know Figure~\ref{fig:ladder} too. A constraint like
$n \le 20$ is not a restriction on the input --- it is a message: \emph{the
intended solution enumerates something}. When you see an unusually small
bound, ask ``what is the intended search space?'' before reaching for anything
cleverer. The clever algorithm that takes an hour to find loses to the
enumeration that takes ten minutes to write and a second to run.
\end{remark}

\begin{example}[WE 2.4 --- USACO Bronze \emph{Milk Pails}: bounds from the
constraints]\label{ex:milkpails}
Farmer John has two pails of sizes $X$ and $Y$, which he may fill completely
any number of times and pour into a third pail of capacity $M$ ($M \le 1000$).
What is the largest amount of milk he can measure into the third pail without
overflowing it?
\par\smallskip
\textbf{Strategy.} What is the search space? Every outcome is ``fill the
$X$ pail $i$ times and the $Y$ pail $j$ times'', giving $iX + jY$ units. The
only question is how far $i$ and $j$ can run --- and the constraint answers
it: $iX + jY \le M \le 1000$, so $i \le M/X$ and $j \le M/Y$ --- worst
case a few hundred thousand pairs. That is a nested loop, not an algorithm: the constraint is
small enough that the search space fits inside it twice over.
\par\smallskip
\begin{solution}
Two loops, both bounded by $M$ itself so no legal combination is clipped:
\lstinputlisting{code/milk_pails.cpp}
On the sample input \texttt{17 25 77}:
\begin{lstlisting}[numbers=none]
$ echo '17 25 77' | ./milk_pails
76
\end{lstlisting}
The inner loop's bound is the line that makes this correct rather than merely
fast: \texttt{i * X + j * Y <= M} clips the search space to candidates that
could still win, and \texttt{best} keeps the largest survivor. Notice what we
did \emph{not} do: no recursion, no pruning rule, no cleverness --- the
constraint gave the bound, the bound gave the loops.
\end{solution}
\end{example}

\section{Enumerating subsets}\label{sec:subsets}
\index{subset!enumerating}

The first enumeration shape you will use weekly: for a set of $n$ elements,
produce each of its $2^n$ \emph{subsets}. Two idioms cover everything.

\emph{Recursion first.} Walk the elements in order; at each one, branch into
``take it'' and ``leave it''. The walk is a binary tree of depth $n$ ---
Figure~\ref{fig:rectree} draws it for $n=3$ --- and every leaf is one subset.

\begin{figure}[ht]\centering
% fig06-rectree.tex -- Figure 2.1: recursion tree for subsets of {1,2,3}.
% Each level decides take/leave for one element; leaves are the 2^3 subsets.
\begin{tikzpicture}[
  font=\footnotesize,
  level distance=8.5mm,
  level 1/.style={sibling distance=48mm},
  level 2/.style={sibling distance=24mm},
  level 3/.style={sibling distance=12mm},
  nd/.style={draw, rounded corners=1pt, inner sep=2pt, fill=graybg},
  lf/.style={draw, rounded corners=1pt, inner sep=2pt, fill=goodbg},
  el/.style={font=\scriptsize, text=slate, midway, above=-1pt},
  every edge/.style={draw, -, thin}
]
\node[nd] {search(1)}
  child { node[nd] {take 1}
    child { node[nd] {take 2}
      child { node[lf] {$\{1,2,3\}$} edge from parent node[el] {3} }
      child { node[lf] {$\{1,2\}$} edge from parent node[el] {} } }
    child { node[nd] {skip 2}
      child { node[lf] {$\{1,3\}$} edge from parent node[el] {3} }
      child { node[lf] {$\{1\}$} edge from parent node[el] {} } }
    edge from parent node[el] {1} }
  child { node[nd] {skip 1}
    child { node[nd] {take 2}
      child { node[lf] {$\{2,3\}$} edge from parent node[el] {3} }
      child { node[lf] {$\{2\}$} edge from parent node[el] {} } }
    child { node[nd] {skip 2}
      child { node[lf] {$\{3\}$} edge from parent node[el] {3} }
      child { node[lf] {$\{\}$} edge from parent node[el] {} } }
    edge from parent node[el] {} };
\end{tikzpicture}

\caption{The recursion tree for subsets of $\{1,2,3\}$. Notice the shape:
$2^n$ leaves at depth $n$, and every internal node is a decision. Drawing
this tree before coding is Section~\ref{sec:sketches}'s whole method.}
\label{fig:rectree}
\end{figure}

\lstinputlisting[caption={\texttt{code/subsets\_rec.cpp}: recursive take/leave.
The skeleton of every backtracking program in this book.},
label={lst:subsetsrec}]{code/subsets_rec.cpp}

Its real output --- the eight subsets in take-first order:

\begin{lstlisting}[numbers=none]
{1, 2, 3}
{1, 2}
{1, 3}
{1}
{2, 3}
{2}
{3}
{}
\end{lstlisting}

\emph{Then bitmasks.}\index{bitmask} A subset of an $n$-element set is also a
number: bit $k$ of a mask says whether element $k$ is in. Then ``all
subsets'' is the loop \texttt{for mask in $[0, 2^n)$} --- the cleanest three
lines in competitive programming, and Figure~\ref{fig:bitmask} shows the
encoding.

\begin{figure}[ht]\centering
% fig07-bitmask.tex -- Figure 2.2: bitmask <-> subset for n = 3.
% One row per mask: the integer, its bits, the subset it encodes.
\begin{tikzpicture}[font=\small, every node/.style={inner sep=0pt}]
\node[draw, text width=0.56\linewidth, inner sep=8pt, font=\footnotesize,
      align=flush left] (tab) at (0,0) {%
  \begin{tabular}{cccl}
  \toprule
  \textbf{mask} & \textbf{bits} $b_2 b_1 b_0$ & \textbf{subset} & \textbf{loop order}\\
  \midrule
  0 & 000 & $\{\}$        & first\\
  1 & 001 & $\{1\}$       &\\
  2 & 010 & $\{2\}$       &\\
  3 & 011 & $\{1,2\}$     &\\
  4 & 100 & $\{3\}$       &\\
  5 & 101 & $\{1,3\}$     &\\
  6 & 110 & $\{2,3\}$     &\\
  7 & 111 & $\{1,2,3\}$   & last\\
  \bottomrule
  \end{tabular}};
\draw[{stealth}-, thick, accent] ($(tab.north east)+(-2.0cm,-0.15)$)
  -- ++(0.5,0.55) node[right, font=\footnotesize, text=accent,
                      text width=3.4cm, align=flush left]
  {bit $k$ set $\iff$ element $k{+}1$ chosen};
\draw[{stealth}-, thick, accent] ($(tab.east)+(0,-0.55)$)
  -- ++(0.5,-0.3) node[right, font=\footnotesize, text=accent,
                      text width=4.0cm, align=flush left]
  {$2^n$ masks cover every subset exactly once};
\end{tikzpicture}

\caption{Mask $\leftrightarrow$ subset for $n=3$. Notice that the loop order
is \emph{numerical}, not the take-first order of recursion: mask $011$ is
$\{1,2\}$, not the next recursive leaf. If your output order matters, check
which order your method produces.}
\label{fig:bitmask}
\end{figure}

\lstinputlisting[caption={\texttt{code/subsets\_mask.cpp}: the same eight
subsets by bitmask. \texttt{mask \& (1 << k)} is the membership test.},
label={lst:subsetsmask}]{code/subsets_mask.cpp}

Real output --- same subsets, different order:

\begin{lstlisting}[numbers=none]
{}
{1}
{2}
{1, 2}
{3}
{1, 3}
{2, 3}
{1, 2, 3}
\end{lstlisting}

\begin{exercise}[easy]\label{exr:subsetcount}
How many subsets does a $20$-element set have? How many masks does the loop
run for $n = 20$ --- and for $n = 64$? What goes wrong in the code of
Listing~\ref{lst:subsetsmask} for large $n$, and where exactly?
\end{exercise}

\begin{example}[WE 2.1 --- CSES \emph{Apple Division}: the enumeration is the
algorithm]\label{ex:apple}
$n$ apples weigh $w_1, \dots, w_n$. Split them into two groups so the weight
difference is as small as possible ($n \le 20$).
\par\smallskip
\textbf{Strategy.} The answer is $\min |S_A - S_B|$ over partitions. Since
$S_A + S_B$ is the fixed total $T$, this is $\min |T - 2S_A|$ --- so we only
need to look at one side: enumerate the subsets that could be group $A$ and
keep the best. The space is $2^n$ with $n \le 20$: at most a million
candidates, comfortably inside the budget by Figure~\ref{fig:ladder}. The
constraint \emph{is} the algorithm.
\par\smallskip
\begin{solution}
\lstinputlisting{code/apple_division.cpp}
On the CSES sample (\texttt{5} then \texttt{3 2 7 4 1}, total $17$ --- best
split puts $3+4+1=8$ in $A$, leaving $|17-16|=1$):
\begin{lstlisting}[numbers=none]
$ echo '5
3 2 7 4 1' | ./apple_division
1
\end{lstlisting}
Two details to copy. First, the loop runs to $2^{n-1}$ not $2^n$: element 0 is
pinned to group $A$, because swapping the two groups changes nothing --- every
answer appears twice in the free search, and the pin deletes the mirror half.
Second, \texttt{best} starts at $T$ (the trivial partition $A = \varnothing$),
so even $n = 1$ works.
\end{solution}
\end{example}

\begin{remark}[Is the symmetry real? Measure it.]
\index{symmetry!halving the tree}
Pinning element~0 halves the search space \emph{if the objective is symmetric}
--- here $|S_A - S_B|$ does not change when groups swap, so the move is free.
\texttt{code/apple\_bench.cpp} times both versions on $n = 22$:
\begin{lstlisting}[numbers=none]
free:      leaves=4194304  time=6.467 ms  answer=1
fixed[0]:  leaves=2097152  time=3.56679 ms  answer=1
\end{lstlisting}
Same answer, half the leaves, roughly half the time. The speedup is modest
because the search was already cheap --- but the \emph{technique} is the
point: before exploiting a symmetry, verify the objective is actually
invariant under it, or you have deleted half your answer space for nothing.
\end{remark}

\section{Enumerating permutations}\label{sec:perms}
\index{permutation!enumerating}

Subsets cost $2^n$; permutations cost $n!$ --- a much harder wall
(Figure~\ref{fig:ladder} shows $n!$ crossing the budget near $n = 11$). When
the constraint allows it, C++ does the walking for you:
\texttt{next\_permutation}\index{next\_permutation@\texttt{next\_permutation}}
advances a sequence to its next arrangement in
\emph{lexicographic order}\index{lexicographic order} --- dictionary order ---
and returns \texttt{false} at the last one.

\begin{tripwire}[The \texttt{sort} contract]
\index{next\_permutation@\texttt{next\_permutation}!must sort first}
\texttt{next\_permutation} only walks \emph{upward}. Called on
\texttt{"bca"}, it produces \texttt{cab}, \texttt{cba} --- then stops, and you
have silently missed four permutations. Sort first, always: from the smallest
arrangement the walk covers all of them. This is the single most common
submission bug in this chapter.
\end{tripwire}

\lstinputlisting[caption={\texttt{code/perms\_demo.cpp}: the idiom up close.
Notice \texttt{"aab"} yields three permutations, not six --- duplicates make
no duplicates.},label={lst:permsdemo}]{code/perms_demo.cpp}

\begin{lstlisting}[numbers=none]
abc:
  abc
  acb
  bac
  bca
  cab
  cba
count: 6
aab:
  aab
  aba
  baa
count: 3
\end{lstlisting}

The second half is worth staring at: repeated elements never produce repeated
permutations, so the count is $n!/\prod(\text{multiplicity}!)$ for free ---
\texttt{"aab"} gives $3!/2! = 3$ distinct strings with zero dedup code.

\begin{example}[WE 2.2 --- CSES \emph{Creating Strings}: apply the
idiom]\label{ex:creating}
Given a string, print every distinct string that can be formed by rearranging
its characters, alphabetically sorted ($n \le 8$).
\par\smallskip
\textbf{Strategy.} $n \le 8$ means at most $8! = 40\,320$ distinct strings ---
inside the budget with room to spare, and the problem's own output size is the
real cost. Sort the input, then let \texttt{next\_permutation} walk; the walk
is already in alphabetical order, so printing as we go satisfies the output
contract for free. Collect-or-print is a choice: we collect to learn the count
first, since the answer's first line must be the total.
\par\smallskip
\begin{solution}
\lstinputlisting{code/creating_strings.cpp}
On \texttt{aabac} ($5!/(3! \cdot 1! \cdot 1!) = 20$ distinct strings), real
output:
\begin{lstlisting}[numbers=none]
20
aaabc
aaacb
aabac
aabca
aacab
aacba
abaac
abaca
abcaa
acaab
acaba
acbaa
baaac
baaca
bacaa
bcaaa
caaab
caaba
cabaa
cbaaa
\end{lstlisting}
The lexicographic order you see is exactly the order the algorithm walks ---
\emph{that's why} sorting first is the contract, not a nicety.
\end{solution}
\end{example}

\begin{exercise}[easy]\label{exr:nextperm}
Write down the permutation \texttt{next\_permutation} produces after
\texttt{3142} --- then after \texttt{4321}. What does each call return?
\end{exercise}

\section{Backtracking}\label{sec:backtracking}
\index{backtracking}

Recursive enumeration has a superpower the loop versions lack: it can see a
dead end early and turn around. \emph{Backtracking} is recursive search with
\emph{pruning}\index{pruning} --- before extending a partial solution, check
whether it can still lead to a valid one, and skip the subtree entirely if
not. The cost model is unchanged (space size $\times$ work per candidate);
pruning attacks the first factor by shrinking the space itself.

A prune is only as good as its rule. Three requirements, in order of
importance: the rule must be \emph{sound} (never kills a branch that contains
an answer --- else the search is wrong, not fast); it should be \emph{cheap}
(a prune that costs more than the subtree it saves is a slowdown wearing a
clever name); and ideally it is \emph{measured} --- which is why the code
below counts.

\begin{example}[WE 2.3 --- CSES \emph{Chessboard and Queens}: pruning with
receipts]\label{ex:queens}
Place eight queens on an $8 \times 8$ board with some squares reserved
(\texttt{*}), so that no two queens share a row, column, or diagonal. Count
the valid placements.
\par\smallskip
\textbf{Strategy.} Place queens row by row --- one per row is forced anyway,
so the state at depth $r$ is a partial placement in rows $0..r-1$. The blind
version tries all $4^8$-style assignments; the smart version checks three
things before recursing: this column already used? this diagonal taken (two
families, indexed $r+c$ and $r-c$)? this square reserved? Any ``yes'' and the
branch dies instantly --- a \emph{prune}, and the program counts them so we
can see the saving rather than guess at it.
\par\smallskip
\begin{solution}
\lstinputlisting{code/queens.cpp}
Run on an empty board, then on the real CSES sample (four reserved squares),
with the diagonal checks switched on and off for comparison:
\begin{lstlisting}[numbers=none]
full checks (col+diag+cells):   nodes=2057   pruned=13664  leaves=92
columns only:                  nodes=109601 pruned=444648 leaves=40320

full checks (col+diag+cells):   nodes=1607   pruned=10730  leaves=65
columns only:                  nodes=68069  pruned=288324 leaves=23520
\end{lstlisting}
Read the counter before the answer. With all checks on, $2057$ nodes produce
$92$ leaves --- and since only fully valid states reach depth 8, every leaf
\emph{is} a solution (CSES sample: $65$). With columns alone, the search still
walks all $8! = 40\,320$ leaf placements --- $53\times$ more nodes visited for
the same information. The diagonal rules did not make the algorithm cleverer;
they deleted 98\% of its work. That is what a pruning rule is for.
\end{solution}
\end{example}

Figure~\ref{fig:queentree} draws the same story on the small board.

\begin{figure}[ht]\centering
% fig08-queentree.tex -- Figure 2.3: unpruned vs pruned search, 4-queens.
% Left: the blind tree fans out 4 ways forever (256 leaves).
% Right: the same search with column+diagonal checks, drawn to scale.
% Counts measured by code/queens.cpp on n=4 (levels 1,4,6,4,2).
\begin{tikzpicture}[font=\scriptsize,
  nd/.style={circle, draw, inner sep=1.2pt, minimum size=9pt, fill=graybg},
  lf/.style={circle, draw, inner sep=1.2pt, minimum size=9pt, fill=goodbg},
  xx/.style={font=\scriptsize, text=amber, inner sep=0pt},
  hd/.style={font=\footnotesize\bfseries},
  lab/.style={font=\scriptsize, text=slate}]

% ---------- left panel: unpruned ----------
\begin{scope}[shift={(0,0)}]
\node[hd] at (0,0.3) {no checks: every branch lives};
\node[nd] (r0) at (0,-0.7) {};
\node[nd] (a) at (-1.5,-1.5) {}; \node[nd] (b) at (-0.5,-1.5) {};
\node[nd] (c) at (0.5,-1.5) {};  \node[nd] (d) at (1.5,-1.5) {};
\draw (r0) -- (a) (r0) -- (b) (r0) -- (c) (r0) -- (d);
\node[lab] at (2.6,-1.5) {row 0: 4};
\foreach \x in {-1.5,-0.5,0.5,1.5} {
  \draw (\x,-1.5) -- ++(-0.28,-0.55) (\x,-1.5) -- ++(-0.09,-0.55)
        (\x,-1.5) -- ++(0.09,-0.55) (\x,-1.5) -- ++(0.28,-0.55);
}
\node[lab] at (2.6,-2.25) {row 1: $4 \times 4 = 16$};
\node[lab] at (2.6,-2.85) {row 2: $4^3 = 64$};
\node[lab] at (2.6,-3.45) {row 3: $4^4 = 256$ leaves};
\end{scope}

% ---------- right panel: pruned, real shape ----------
\begin{scope}[shift={(6.6,0)}, scale=1.05]
\node[hd] at (0,0.3) {column + diagonal checks};
\node[nd] (q) at (0,-0.7) {};
% row 0 children: columns 0..3
\node[nd] (p0) at (-2.4,-1.5) {0}; \node[nd] (p1) at (-0.8,-1.5) {1};
\node[nd] (p2) at (0.8,-1.5) {2};  \node[nd] (p3) at (2.4,-1.5) {3};
\draw (q) -- (p0) (q) -- (p1) (q) -- (p2) (q) -- (p3);
% row 1 survivors
\node[nd] (p02) at (-2.85,-2.4) {2}; \node[nd] (p03) at (-2.0,-2.4) {3};
\node[nd] (p13) at (-0.8,-2.4) {3};
\node[nd] (p20) at (0.8,-2.4) {0};
\node[nd] (p30) at (2.0,-2.4) {0}; \node[nd] (p31) at (2.8,-2.4) {1};
\draw (p0) -- (p02) (p0) -- (p03) (p1) -- (p13)
      (p2) -- (p20) (p3) -- (p30) (p3) -- (p31);
% pruned row-1 stubs (two for p0, three for p1, three for p2, two for p3)
\node[xx] at (-2.05,-2.1) {\sf x}; \node[xx] at (-1.65,-2.1) {\sf x};
\node[xx] at (-1.0,-2.1) {\sf x}; \node[xx] at (-0.55,-2.1) {\sf x};
\node[xx] at (0.45,-2.1) {\sf x};  \node[xx] at (1.05,-2.1) {\sf x};
\node[xx] at (2.6,-2.1) {\sf x};   \node[xx] at (3.15,-2.1) {\sf x};
\node[lab] at (4.6,-2.1) {row 1: 6 of 16};
% row 2 survivors
\node[nd] (r21) at (-2.0,-3.3) {1};
\node[nd] (r20) at (-0.8,-3.3) {0};
\node[nd] (r23) at (0.8,-3.3) {3};
\node[nd] (r22) at (2.0,-3.3) {2};
\draw (p03) -- (r21) (p13) -- (r20) (p20) -- (r23) (p30) -- (r22);
\node[xx] at (-2.85,-2.95) {\sf x x x x};
\node[xx] at (0.75,-2.95) {\sf x x x};
\node[xx] at (2.8,-2.95) {\sf x x x x};
\node[lab] at (4.6,-3.05) {row 2: 4 of 64};
% row 3: two solutions
\node[lf] (s1) at (-0.8,-4.2) {2};
\node[lf] (s2) at (0.8,-4.2) {1};
\draw (r20) -- (s1) (r23) -- (s2);
\node[xx] at (-2.0,-3.85) {\sf x x x x};
\node[xx] at (2.0,-3.85) {\sf x x x x};
\node[lab] at (4.6,-3.95) {row 3: 2 solutions};
\end{scope}
\end{tikzpicture}

\caption{4-queens, unpruned vs pruned (real shape, measured). Left: the blind
tree fans out to $4^4 = 256$ leaves. Right: with column + diagonal checks, the
row-1 level keeps 6 of 16 children, row 2 keeps 4 of 64, and exactly 2
solutions survive. Every \textsf{x} is a subtree that was never entered ---
the saving is the size of what is missing, not what is present.}
\label{fig:queentree}
\end{figure}

\begin{exercise}[medium]\label{exr:prunefire}
In \texttt{code/queens.cpp}, the three prune rules live inside one boolean
expression. Which exact line of \texttt{search()} fires when a candidate
square shares a column with a placed queen? If you removed the
\texttt{useDiag} term but kept \texttt{useCell}, would \texttt{leaves} still
equal the answer? Why or why not?
\end{exercise}

\section{Search-space sketches}\label{sec:sketches}
\index{search space!sketching}

Every problem in this chapter was planned the same way: draw the tree before
coding it. The sketch costs a minute and answers three questions an IDE
cannot: How big is the space? What does a node look like? Where can a check
fire?

The method is mechanical. First, \emph{name the state}: the minimal
information a recursion call needs (for subsets, ``depth $k$ + what's been
chosen''; for queens, ``which columns are taken''). Second, \emph{draw two
levels of the tree} on a small instance ($n = 3$ or $4$) --- if the leaves are
not obviously the candidates, the state is wrong. Third, \emph{count the
leaves and the internal nodes}: leaves tell you the space size, internals tell
you the real work, and their ratio tells you how much a per-node prune could
save. Figure~\ref{fig:rectree} and Figure~\ref{fig:queentree} are exactly
these sketches, left in the book on purpose --- they are not illustrations of
the code, they are the \emph{design document} the code was written from.

\begin{exercise}[medium]\label{exr:trace}
Draw the full recursion tree for \texttt{subsets\_rec.cpp} on $n = 4$, marking
take/leave on every edge as in Figure~\ref{fig:rectree}. How many internal
nodes and leaves does it have? Now mark the tree for the \emph{fixed-element}
version of WE~2.1's search: which half disappears, and which node does the
pin live at?
\end{exercise}

\section{Common mistakes}
\index{common mistakes}

\begin{itemize}
\item \textbf{\texttt{1 << n} on the wrong side of 31.} \texttt{1} is an
\texttt{int}; shifting it into bit 31 is undefined and bit 32+ is garbage.
For $n \ge 31$ write \texttt{1LL << n} --- and for $n \ge 64$ admit that a
single mask cannot hold the subset at all (use \texttt{vector<bool>} or the
recursive version, which never builds a mask).
\item \textbf{\texttt{next\_permutation} on unsorted input.} The walk goes
upward only: from \texttt{"cba"} it ends immediately and you enumerate one
permutation, silently. Sort first --- then duplicates need no other handling.
\item \textbf{Marking without unmarking.} Backtracking's bookkeeping is
borrowed state: \texttt{col[c] = true} before the recursive call,
\texttt{col[c] = false} after. Forget the second line and the next sibling
inherits phantom conflicts --- every subsequent branch dies for no reason and
the count is quietly wrong. The mark/unmark pair is the price of the prune.
\item \textbf{Halving on a non-symmetric objective.} The WE~2.1 pin worked
because $|S_A - S_B|$ is unchanged by swapping the groups. If the objective
distinguishes the sides --- ``first group's sum must not exceed the
second's'' --- the same pin deletes exactly the answers you needed. Verify
the invariance on paper before deleting half the tree.
\item \textbf{Counting leaves as the cost.} The tree's work happens at
\emph{internal} nodes too: the subset tree of Figure~\ref{fig:rectree} has
$2^n$ leaves but $2^{n+1} - 1$ nodes total. When estimating a search, count
nodes, not answers --- and when claiming a prune helps, measure both counters
as WE~2.3 does, because a rule that prunes the wrong branch is worse than
none.
\end{itemize}

\section{Chapter summary}
\begin{table}[ht]\centering
\begin{tabularx}{\linewidth}{@{}lYY@{}}
\toprule
\tablehead\thead{Tool} & \thead{Use when} & \thead{Cost / first check when wrong}\\
\midrule
Nested loops & space is a small product of bounds ($\le 10^6$ pairs) & bound clips a legal case: check $iX + jY \le M$ style\\
Recursive subsets & you want take/leave order or per-node work & $2^n$ leaves, $2^{n+1}-1$ nodes; depth $= n$\\
Bitmask \texttt{for mask} & you want three lines and numerical order & \texttt{1 << n} overflows past $n = 30$\\
\texttt{next\_permutation} & all $n!$ orders needed, $n \le 11$ & \texttt{sort} first or the walk is incomplete\\
Backtracking & dead branches are detectable early & prune must be sound --- never kills an answer\\
Symmetry pin & objective invariant under some swap & verify invariance, then halve --- measure it\\
\bottomrule
\end{tabularx}
\caption{Chapter 2 in one page. The meta-rule sits above the table: estimate
the space before you enumerate it --- size times work per candidate decides
everything.}
\label{tab:ch2summary}
\end{table}

\section{Additional exercises}
\index{exercises}
Named archive problems are cited by name only; look them up on
\texttt{cses.fi} and \texttt{usaco.org}.

\begin{exercise}[easy] CSES \emph{Tower of Hanoi}: print the optimal move
sequence for $n$ disks. [implement --- this is a recursive walk that emits
its own search tree]
\end{exercise}
\begin{exercise}[easy] CSES \emph{Gray Code}: list all $2^n$ bit strings so
adjacent strings differ in one bit. [implement --- one recursive flip order
produces the whole list]
\end{exercise}
\begin{exercise}[easy] USACO Bronze \emph{Diamond Collector}: choose the two
largest diamonds whose sizes differ by at most $k$. Trains bounded nested
loops and the ``sort then think'' reflex.
\end{exercise}
\begin{exercise}[easy] USACO Bronze \emph{Daisy Chains}: count plants whose
height equals the average height of some contiguous interval. Nested loops
over interval endpoints --- the space is $O(n^2)$, the check is $O(n)$.
\end{exercise}
\begin{exercise}[easy] USACO Bronze \emph{Counting Liars}: find the earliest
true statement in a list of claims about a hidden position. Sort the claims
and walk --- complete search over $n$ candidates.
\end{exercise}
\begin{exercise}[medium] CSES \emph{Raab Game I}: construct a game history
with given scores, or report impossible. Search meets construction.
\end{exercise}
\begin{exercise}[medium] CSES \emph{Mex Grid Construction}: fill an $n
\times n$ grid so each cell's mex condition holds. Enumeration under a local
constraint --- a backtracking-shaped construction.
\end{exercise}
\begin{exercise}[medium] USACO Bronze \emph{Cow Gymnastics}: from $K$
practice rankings, count pairs of cows whose order never changes. Triple
nested loops bounded by the constraints.
\end{exercise}
\begin{exercise}[medium] USACO Bronze \emph{Triangles}: sum the areas of all
axis-aligned right triangles in a point set. Enumeration over $\binom{n}{2}$
or $\binom{n}{3}$ --- choose which.
\end{exercise}
\begin{exercise}[hard] CSES \emph{Digit Queries}: find the $k$-th digit of
the concatenated sequence $1, 2, 3, \dots$ --- the search space is far past
enumeration, so derive the block containing $k$ instead. [derive --- a lesson
in when \emph{not} to enumerate]
\end{exercise}
\begin{exercise}[hard] CSES \emph{Grid Path Description}: count paths
visiting every square of a $7 \times 7$ grid along a given pattern. Pure
backtracking --- and it only fits inside the budget once you add the pruning
rules WE~2.3 taught. [implement]
\end{exercise}

\chapternotes
The complete-search chapter of Laaksonen's handbook~\cite{laaksonen2024}
covers the same three idioms in the same order --- generating subsets both
ways, permutations, then backtracking with pruning --- and this chapter's
worked examples are drawn from the CSES introductory set it
accompanies~\cite{cses}. Sannemo's treatment adds the ``search space size
times per-candidate work'' cost model used in Section~\ref{sec:exhaustive}
and the ``small constraint means intended enumeration'' reading of
Remark~\ref{rem:smalln}~\cite{sannemo2025}. The USACO Guide's Bronze modules
on complete search and recursion supplied the nested-loop framing of WE~2.4
and several of the additional-exercise choices~\cite{usacoguide}; the
backtracking instrumentation (count the nodes, count the prunes) follows the
measure-first debugging style of cp-algorithms' backtracking
articles~\cite{cpalgorithms}. All timings are real runs on this machine
(Apple clang, \texttt{-O2}) and will differ modestly elsewhere --- the
ratios, not the numbers, carry the argument. Halim's text~\cite{halimcp4}
organises the same material around UVa problem families --- cited here, not
consulted for content.

